\noindent\textit{Sửa từ đề cũ số 4. Biên soạn: Trang.}

\begin{bt}[2,0 điểm]
Thực hiện phép tính:
\begin{listEX}[1]
\item $\left(-\dfrac{3}{7}+\dfrac{4}{11}\right):\dfrac{7}{11}+\left(-\dfrac{4}{7}+\dfrac{7}{11}\right):\dfrac{7}{11}$
\item $\sqrt{\dfrac{25}{36}}\cdot 1{,}5+\left(-\dfrac{1}{4}\right)^{2}:\left|-0{,}25\right|-\left(\dfrac{7}{8}\right)^{0}$
\item $\dfrac{10^{4}\cdot 81-16\cdot 15^{2}}{4^{4}\cdot 675}$
\item $\dfrac{6^{5}\cdot 12^{3}}{8^{4}\cdot 27^{2}}$
\end{listEX}
\loigiai{
a) $0$.\quad
b) $\dfrac{1}{2}$.\quad
c) $\dfrac{14}{3}$.\quad
d) $\dfrac{9}{2}$.
}
\end{bt}

\begin{bt}[2,0 điểm]
Tìm $x$, biết:
\begin{listEX}[2]
\item $\dfrac{11}{2}+\dfrac{4}{5}x=\dfrac{5}{6}$
\item $\left|x-\dfrac{1}{2}\right|=\dfrac{3}{4}$
\item $4\left(\dfrac{2x-5}{6}\right)^{3}=-\dfrac{32}{27}$
\item $5^{x+1}+5^{x}=150$
\end{listEX}
\loigiai{
a) $x=-\dfrac{35}{6}$.\quad
b) $x=\dfrac{5}{4}$ hoặc $x=-\dfrac{1}{4}$.\quad
c) $x=\dfrac{1}{2}$.\quad
d) $x=2$.
}
\end{bt}

\begin{bt}[1,0 điểm]
Một cửa hàng nhân dịp khai trương giảm giá $15\%$ so với giá niêm yết cho tất cả các mặt hàng.
\begin{listEX}[1]
\item Biết giá niêm yết của một chiếc tivi là $14\,000\,000$ đồng. Tính giá bán của tivi sau khi giảm.
\item Anh Bình mua một tivi và một máy lạnh, phải trả tất cả $22\,100\,000$ đồng. Tính giá niêm yết của máy lạnh.
\end{listEX}
\loigiai{
a) $11\,900\,000$ đồng.\\
b) $12\,000\,000$ đồng. Giá niêm yết của cả hai món là $22\,100\,000:0{,}85=26\,000\,000$ đồng.
}
\end{bt}

\begin{bt}[1,0 điểm]
Một xưởng dự định hoàn thành một lô hàng trong $12$ ngày. Nếu mỗi ngày xưởng làm thêm $5$ sản phẩm thì hoàn thành sớm $2$ ngày và làm được thêm $4$ sản phẩm so với số sản phẩm dự định của cả lô. Hỏi theo kế hoạch, mỗi ngày xưởng làm bao nhiêu sản phẩm và cả lô có bao nhiêu sản phẩm?
\loigiai{
Mỗi ngày $23$ sản phẩm; cả lô $276$ sản phẩm.\\
Gọi $x$ là số sản phẩm làm trong một ngày theo kế hoạch ($x>0$). Ta có $10(x+5)=12x+4$.
}
\end{bt}

\begin{bt}[1,5 điểm]
Cho đường thẳng $a$ vuông góc với đường thẳng $b$ tại $B$ và vuông góc với đường thẳng $c$ tại $C$. Đường thẳng $d$ cắt $b$ tại $M$, cắt $c$ tại $N$ và tạo với $b$ một góc $120^{\circ}$.
\begin{listEX}[1]
\item Chứng minh $b\parallel c$.
\item Tính góc nhọn tạo bởi $d$ và $c$.
\item Vẽ tia $Nt$ là tia phân giác của góc nhọn tạo bởi $d$ và $c$. Tính góc giữa tia $Nt$ và đường thẳng $c$.
\end{listEX}
\begin{center}
\begin{tikzpicture}[scale=1.05,line cap=round]
	\draw (0,2.2) -- (5.2,2.2) node[right] {$b$};
	\draw (0,0) -- (5.2,0) node[right] {$c$};
	\coordinate (B) at (4.1,2.2);
	\coordinate (C) at (4.1,0);
	\coordinate (M) at (1.15,2.2);
	\coordinate (N) at (2.45,0);
	\draw (4.1,2.9) -- (4.1,-0.55) node[below] {$a$};
	\draw (M) -- (N);
	\draw (N) -- ($(N)+(0.55,-0.45)$) node[below] {$d$};
	\coordinate (ML) at (0.2,2.2);
	\node[above right] at (B) {$B$};
	\node[below right] at (C) {$C$};
	\node[above] at (M) {$M$};
	\node[below left] at (N) {$N$};
	\pic[draw,angle radius=3mm] {right angle = M--B--C};
	\pic[draw,angle radius=3mm] {right angle = N--C--B};
	\pic[draw,angle radius=5mm,angle eccentricity=1.55,"$120^{\circ}$"] {angle = ML--M--N};
\end{tikzpicture}
\end{center}
\loigiai{
a) $b\perp a$ và $c\perp a$ nên $b\parallel c$.\\
b) $60^{\circ}$. Vì $b\parallel c$ nên góc tương ứng bằng $120^{\circ}$, góc nhọn kề bù với góc đó bằng $60^{\circ}$.\\
c) $30^{\circ}$.
}
\end{bt}

\begin{bt}[2,0 điểm]
Cho tam giác $ABC$ vuông tại $A$ và $\widehat{ABC}=40^{\circ}$. Điểm $D$ thuộc cạnh $BC$ sao cho $BD=AB$. Trên tia đối của tia $DA$ lấy điểm $E$ sao cho $DE=DC$.
\begin{listEX}[1]
\item Tính $\widehat{ACB}$. Chứng minh $\triangle ABD$ cân và tính các góc của tam giác đó.
\item Tính $\widehat{CAD}$ và $\widehat{ADC}$.
\item Chứng minh $\triangle CDE$ cân và tính số đo hai góc ở đáy.
\end{listEX}
\loigiai{
a) $\widehat{ACB}=50^{\circ}$. Tam giác $ABD$ cân tại $B$, $\widehat{ABD}=40^{\circ}$, $\widehat{BAD}=\widehat{ADB}=70^{\circ}$.\\
b) $\widehat{CAD}=20^{\circ}$, $\widehat{ADC}=110^{\circ}$.\\
c) $CD=DE$ nên $\triangle CDE$ cân tại $D$. Góc ở đỉnh $D$ bằng $70^{\circ}$, hai góc ở đáy bằng $55^{\circ}$.
}
\end{bt}

\begin{bt}[0,5 điểm]
Tính giá trị của biểu thức
\[
M=\dfrac{2^{2}}{1\cdot 3}\cdot\dfrac{3^{2}}{2\cdot 4}\cdot\dfrac{4^{2}}{3\cdot 5}\cdots\dfrac{20^{2}}{19\cdot 21}.
\]
\loigiai{
$M=\dfrac{40}{21}$.\\
\textit{Hướng dẫn.} Rút gọn chéo, còn $\dfrac{2}{1}\cdot\dfrac{20}{21}$.
}
\end{bt}
